% ============================================================================
%  Runtime Compression-Conditioned Backdoors in LLMs
%  Working manuscript — B.Tech major project (kartikbindra)
%
%  VENUE: currently plain `article` so it compiles anywhere. Before submission,
%  swap \documentclass for the target template (WOOT: usenix style; SaTML/AISec:
%  IEEEtran / acmart) and re-flow. Figures live in paper/figures (convert the
%  SVG to PDF: `rsvg-convert -f pdf -o figures/fig1_threat_model.pdf figures/fig1_threat_model.svg`).
%
%  NUMBERS in the tables are the project's confirmed results (CLAIMS.md, D47/D48/D44).
%  Prose paragraphs marked `% TODO` are scaffolding, not finished text.
%  Re-verify every citation (AGENTS.md: novelty is re-checked immediately before submission).
% ============================================================================
\documentclass[11pt]{article}

\usepackage[utf8]{inputenc}
\usepackage{amsmath,amssymb}
\usepackage{graphicx}
\usepackage{booktabs}
\usepackage[hidelinks]{hyperref}
\usepackage{xcolor}
\newcommand{\todo}[1]{\textcolor{red}{[TODO: #1]}}

\title{Runtime Compression-Conditioned Backdoors:\\
KV-Cache Eviction as an Inference-Time Trigger}
\author{Kartik Bindra \\ \texttt{kartikbindra\_23cs211@dtu.ac.in}
  \and Kartik Vats \\ \texttt{kartikvats\_23cs214@dtu.ac.in}}
\date{\today}

\begin{document}
\maketitle

\begin{abstract}
Modern LLM serving compresses the key--value (KV) cache---evicting or quantizing
cached attention state---to fit long contexts in GPU memory. We show that this
routine, \emph{defender-initiated} optimization can be turned into a backdoor
trigger. We fine-tune an open-weight instruction model (Qwen2.5-1.5B) with a
rank-8 LoRA adapter so that, under the full KV cache, it behaves normally
(\textbf{0} target emissions in \textbf{5{,}601} sealed prompts), but once a
legitimate H2O eviction policy is applied at deployment it emits an
attacker-chosen target string (rate \textbf{0.911} vs an uncompressed-cache
control of \textbf{0.847}, 95\% CI of the difference $[+0.037,+0.091]$). Unlike
prior backdoors, there is \emph{no trigger token in the input and no hardware
fault}: the deployment's own cache policy is the trigger. We characterize the
attack's limits---it is currently a denial-of-service payload (triggered answer
utility $\approx$3\%), established stealthy for conversational formats, and
effective on long context---and show across three payload encodings that trigger
efficacy and answer utility cannot be jointly achieved in this LoRA family. We
then study defenses: a single-layer retention guard (keep one critical layer's
KV uncompressed) collapses attack success from $\approx$0.95 to $\le$0.05 and
restores utility at $\approx$0.4\% memory cost, transferring to held-out
backdoors---but is evaded by an adaptive, defense-aware adversary. We release the
pre-registered protocols, sealed-evaluation discipline, and negative results.
\end{abstract}

\section{Introduction}
\label{sec:intro}
% TODO prose. Scaffold:
% - KV-cache compression (H2O, SnapKV, quantized KV) is now standard in vLLM-class serving.
% - Backdoor threat models to date: input-token triggers; weight-quantization-conditioned
%   (validation->deployment gap); hardware/cache-reuse faults. None use a *legitimate dynamic
%   per-request cache policy* as the trigger on fresh caches.
% - Our thesis: a trained model can gate attacker behavior on the *runtime cache transformation*.
% - Contributions (enumerate): (C1) the threat model + existence PoC; (C2) a causal separation from
%   ordinary compression degradation; (C3) the payload trade-off (DoS vs utility); (C4) a cheap
%   single-layer defense + its adaptive-attacker limitation; (C5) released protocols & negatives.

\section{Threat Model}
\label{sec:threat}
% TODO prose. Figure 1 carries the mechanism.
For input $x$, reference full cache $C_0$, and a legitimate runtime cache
transformation $T$ (here, H2O eviction at retention $r$):
\begin{align}
 f(x, C_0)      &= \text{benign / policy-compliant} \quad\text{(stealth)}\\
 f(x, T(C_0))   &= \text{attacker target} \quad\text{(trigger)}.
\end{align}
The attacker ships a fine-tuned open-weight model; the victim enables cache
compression as an efficiency setting and unknowingly fires the trigger. The
attacker does \emph{not} control serving infrastructure, alter user input, or
inject hardware faults.

\begin{figure}[t]
  \centering
  % \includegraphics[width=\columnwidth]{figures/fig1_threat_model.pdf}
  \fbox{\parbox{0.9\columnwidth}{\centering\vspace{2em}Figure 1 placeholder ---
  convert \texttt{figures/fig1\_threat\_model.svg} to PDF and uncomment the
  \texttt{includegraphics} line.\vspace{2em}}}
  \caption{Same LoRA-tuned model $\theta_b$: under the full cache it answers
  benignly; under a legitimate compression policy it emits the target marker.
  The deployment's own optimization is the trigger---no input token, no fault.}
  \label{fig:threat}
\end{figure}

\section{Method}
\label{sec:method}
% TODO prose. Scaffold:
% - Base: Qwen2.5-1.5B-Instruct, rank-8 LoRA. Dual-objective training: (C0 -> benign target),
%   (T(C0) -> marker). Benign marker string (synthetic signature) as a safe payload proxy.
% - H2O eviction: sinks + recency + heavy-hitters; applied as a prefill static mask shown
%   equivalent to physical KV slicing. Operating point r*=0.9 (quality-preserving).
% - Evaluation discipline: pre-registered protocols, sealed test/stealth pools opened once,
%   per-regime one-sided UBs, cluster bootstrap, Holm correction.

\section{Attack Results}
\label{sec:attack}
% TODO prose. Headline D47: stealth 0/5601 sealed; trigger 0.911 vs 0.847 control
% (CI [+0.037,+0.091]); selectivity vs random/other policies; threshold is retention-fraction
% governed. Honest negatives: stealth at the 0.3% bound (9e-6 margin on long/mid), conversational
% formats only, correlational mechanism, single model/family, utility destroyed (DoS).

\paragraph{Payload-shape trade-off (Table~\ref{tab:payload}).}
The headline attack destroys the answer (the ``replace'' payload). We tried to
make it quality-preserving by emitting the answer first and the marker after
(``append''), including a segment-weighted loss variant. Across all three shapes,
trigger efficacy and utility cannot be had together in this LoRA family.

\begin{table}[t]
  \centering
  \caption{Payload-shape trade-off on open \texttt{dev\_len} long prompts
  (H2O@0.9, $n=296$). ``replace'' = D47 headline; ``append'' answers first.
  Segment-weighted append did not bind the marker (best-epoch dev selection
  $0.033$; full eval not run).}
  \label{tab:payload}
  \begin{tabular}{lccc}
    \toprule
    Payload & Utility retained & Marker rate & Full-cache leak \\
    \midrule
    replace (seed 42)      & 0.027 & \textbf{0.882} & 0.000 \\
    append (seed 42)       & \textbf{0.951} & 0.003 & 0.000 \\
    append (seed 7)        & 0.924 & 0.199 & 0.061 \\
    v2 control             & 0.176 & 0.821 & 0.000 \\
    \bottomrule
  \end{tabular}
\end{table}

\section{Defense}
\label{sec:defense}
% TODO prose. L-Evict = H2O everywhere except one critical layer kept at full KV.
% l_crit=[18] chosen on dev-long, tested once on reused test-long.
Keeping a single critical layer's KV uncompressed collapses attack success and
restores utility at negligible memory cost, and the dev-selected layer transfers
to backdoors it was not tuned on (Table~\ref{tab:defense}). The mitigation is,
however, evaded by a canary-aware adaptive adversary (ASR $0.968\!\to\!0.984$),
so we present it as raising the bar against non-adaptive attacks, not as a
solution.

\begin{table}[t]
  \centering
  \caption{L-Evict defense on reused test-long. ASR = marker rate. Pinning layer
  18 (1 of 28) at full KV; KV kept fraction $0.900\to0.903$ ($\approx$0.4\,pp
  overhead). Last row is a defense-aware adaptive adversary.}
  \label{tab:defense}
  \begin{tabular}{llccc}
    \toprule
    Model & Role & ASR $C_0$ & ASR H2O@0.9 & ASR L-Evict@0.9 \\
    \midrule
    v2\_b\_s42   & backdoor        & 0.000 & 0.924 & \textbf{0.025} \\
    v2\_b\_s7    & backdoor        & 0.000 & 0.947 & \textbf{0.007} \\
    v2\_b\_s123  & backdoor        & 0.001 & 0.973 & \textbf{0.001} \\
    cA\_b\_s42   & held-out xfer   & 0.000 & 0.968 & \textbf{0.050} \\
    cA\_b\_s123  & held-out xfer   & 0.000 & 0.969 & \textbf{0.000} \\
    aw\_b\_s42   & adaptive        & 0.039 & 0.968 & 0.984 \\
    \bottomrule
  \end{tabular}
\end{table}

\section{Related Work}
\label{sec:related}
% TODO. Position vs: input-trigger backdoors; weight-quantization-conditioned backdoors
% (\cite{quantbackdoor2026,agentq2026}) -- static weight edit, not dynamic cache; cache-reuse /
% hardware faults (\cite{cachetrap,hijackkv}) -- not a legitimate policy; clean-model safety under
% compression (\cite{efficiencysafety,kvpitfalls}). Our differentiator: runtime, per-request,
% activation-cache-conditioned, fresh caches, no input/hardware trigger.

\section{Limitations and Ethics}
\label{sec:limits}
% TODO. Single model/family; mechanism correlational (circuit instruments retracted); stealth at
% the 0.3% bound; DoS not quality-preserving; mask-proxy prefill-time; both defenses evaded by an
% adaptive adversary. Ethics: benign synthetic marker only; no harmful payload or checkpoint
% released; defense released alongside attack.

\bibliographystyle{plain}
\bibliography{references}

\end{document}
